\pdfoutput=1
% 6+6 restructure (research lab draft, 2026-11-24). The current paper (main.tex) stays buildable; the lead engineer
% integrates. Every number below resolves to claims/claims.yaml (section LAB). Items marked \pend{...} are not
% results and must not be stated until their result file exists.
\documentclass[11pt]{article}
\usepackage{emnlp2023}
\usepackage{times}
\usepackage{latexsym}
\usepackage[T1]{fontenc}
\usepackage[utf8]{inputenc}
\usepackage{microtype}
\usepackage{booktabs}
\usepackage{graphicx}
\usepackage{amsmath,amssymb,amsthm}
\newtheorem{theorem}{Theorem}
\newtheorem{corollary}{Corollary}
\newcommand{\RV}{R$_V$}
\newcommand{\pend}[1]{\textbf{[pending: #1]}}
\newcommand{\sgn}{\operatorname{sign}}

\title{Group Normalisation Turns the Value of Checking into a Vote:\\Cost-Blind Verification in RL-Trained Tool Agents}

\author{Atharv Sharma \\
  BITS Pilani, Goa Campus \\
  \texttt{atharvgaur591@gmail.com}}

\begin{document}
\maketitle

\begin{abstract}
A tool agent whose writes can fail silently must decide when to pay for a read that checks the authoritative state.
The right rule is value of information: verify when the probability that checking changes the outcome, times what
the outcome is worth, exceeds the cost. We show that GRPO's group-standardised advantage does not learn this rule
when the rollouts of a group share a hidden execution state, as they do whenever a task's fault plan is fixed per
prompt. For a decision whose payoff is deterministic given the hidden state, GRPO's expected update is exactly the
probability-weighted \emph{majority sign} of the action's effect across states, while Dr.~GRPO and RLOO estimate its
expected effect (Theorem~1, exact at finite group size). GRPO therefore verifies iff checking changes the outcome in
more than half of the states, whatever the cost, the task value or the failure penalty. In an exactly solvable
costed-verification testbed, the sign of this exact expected update at the warm start predicts the decision GRPO
learns in 494 of 495 cells; GRPO decides wrongly in 64\% of them, Dr.~GRPO in 6\%. A causal intervention that keeps
the hidden-state-conditional scale but removes sample sharing reproduces the failure, and a scale of the same kind
computed without the hidden state removes it. In Tiny Aya Global trained on ProcureSim, the group-std arm stops
reading in both development seeds and the training-log decomposition shows the predicted cancellation. A natural
remedy, counterfactual decision-value credit from paired forks of the simulator state, is a known estimator family
(TRPO's vine, COMA's counterfactual advantage, CRR); at matched episode budgets it does not beat the strongest simple
baseline, with or without adaptive branch allocation. The practical conclusion is a one-line change: do not divide
by a scale estimated from samples that share a hidden state.
\end{abstract}

\section{Introduction}

A payment API returns \texttt{\{"status": "ok"\}}; the payment may have been stored, stored in part, stored twice, or
not at all. An agent can read the payment record to find out, at a cost in latency, money or rate limit. Whether to
read is a value-of-information (VOI) decision: read iff $p\,s\,(V+L_f) > c$, where $p$ is the probability of a silent
fault, $s$ the probability that a detected fault is repaired, $V$ the task value, $L_f$ the failure penalty and $c$
the cost of the read.

Group-relative policy optimisation (GRPO) standardises each rollout's return within its group,
$\hat A_i=(R_i-\bar R_g)/(\sigma_g+\epsilon)$. When a tool-agent task is posed with a fixed hidden fault plan, every
rollout in a group shares that hidden state. We show that this combination removes the cost from the learned rule
entirely.

\paragraph{Contributions.}
(i) \textbf{A finite-group theorem} (\S\ref{sec:mech}). With a hidden state shared within the group and payoffs
deterministic given the state, GRPO's expected update on a binary decision is
$(G{-}1)\,m_G(\rho)\sum_z\pi(z)\sgn\beta_z$, where $\beta_z$ is the action's effect in state $z$. Dr.~GRPO's is
$(G{-}1)\rho(1{-}\rho)\,\mathrm{VOI}$ exactly. The result extends the group-std identity of
\citet{bay2026identity} from binary rewards to arbitrary stakes, and it applies the hidden-state version of the
$1/\sigma$ weighting of \citet{bereket2025uncalibrated} to a costed information action. We state when the effect
weakens and when it vanishes.
(ii) \textbf{Exact confirmation} in a solvable costed-verification testbed with 18 estimators (\S\ref{sec:exp}),
including a causal intervention on the normaliser and an exact gradient bias/variance study.
(iii) \textbf{Observation in a trained language agent}: a reproduction of the training-log decomposition, and
development runs in which the group-std arm stops reading.
(iv) \textbf{A negative result for the obvious remedy.} Paired-fork decision-value credit is a known family, and it
does not pay for its forks at matched budget, even with adaptive branch allocation.

\section{Problem and environment}\label{sec:setup}

\paragraph{POMDP.} A task $x$ has a hidden execution state $z\sim\pi_0(\cdot\mid x)$: which writes silently failed,
were applied partially, duplicated or timed out. The agent sees histories $h_t$ of tool calls and replies. A write's
\emph{ok} acknowledgement is identical in every $z$. After an acknowledgement the agent chooses $r\in\{0,1\}$
(continue or verify), with $\rho=\pi_\theta(r{=}1\mid h)$. The return is
$R=V\cdot\mathrm{succ}-c\,n_{\mathrm{reads}}-L_f(1-\mathrm{succ})-D\,n_{\mathrm{dup}}$.
The deterministic SQLite verifier, not the agent, decides $\mathrm{succ}$.

\paragraph{ProcureSim and the costed-verification pool.} ProcureSim (Appendix~\ref{app2:env}) processes purchase
orders through order, receive, pay and post tools. Its faults (phantom, partial, lost acknowledgement, duplicate
delivery, never-commit) return byte-identical acknowledgements. In the costed pool (\texttt{cw}), each read costs
$c=0.05$ and the reward is $R_V-c\,n_{\mathrm{reads}}$. Training uses Tiny Aya Global (3.35B) with LoRA rank 16,
warm-started by SFT, then 200 steps of TRL GRPO (4 prompts $\times$ 8 rollouts). The task's fault plan is fixed
per prompt, so the 8 rollouts of a group share $z$ (``condition A'').

\paragraph{Exact testbed.} A CPU simulator with the same structure: writes with fault rate $p$, repair rate $s$,
false alarms, visible timeouts with duplicate risk, and a never-faulted write. The optimum is computed exactly by
expectimax over the hidden fault plan, and so is every policy's value, so regret $J^*-J$ is exact. The policy is a
softmax over observable features only; tests check that the hidden state never enters an actor or critic input.

\section{The mechanism: a stake-free vote}\label{sec:mech}

Let $k\sim\mathrm{Bin}(G,\rho)$ be the number of verifying rollouts in a group and
$m_G(\rho)=\mathbb E\big[\sqrt{k(G-k)/(G(G-1))}\,;\,0<k<G\big]$.

\begin{theorem}\label{thm:vote}
Suppose all $G$ rollouts of a group share $z$, and $R_i=\alpha_z+\beta_z r_i$ (deterministic within the state).
Then as $\epsilon\to0$ the expected update on the verify logit satisfies
$U_{\mathrm{GRPO}}=(G{-}1)\,m_G(\rho)\sum_z\pi_0(z)\sgn(\beta_z)$,
$U_{\mathrm{Dr}}=(G{-}1)\rho(1{-}\rho)\sum_z\pi_0(z)\beta_z$ and
$U_{\mathrm{RLOO}}=G\rho(1{-}\rho)\sum_z\pi_0(z)\beta_z$.
The last two need only i.i.d.\ rollouts given the state(s), so they hold for stochastic payoffs and for i.i.d.\
states too.
\end{theorem}
\noindent\emph{Proof sketch.} Inside state $z$, $\hat A_i=\sgn(\beta_z)(r_i-\bar r)/\hat\sigma_r$. With the score
$r_i-\rho$ and $\sum_i(r_i-\bar r)=0$, the group contributes $\sgn(\beta_z)(G-1)\hat\sigma_r$. Taking expectations
over $k$ gives the result. The magnitude $|\beta_z|$ cancels; without the std it survives
(Appendix~\ref{app2:proofs}). $\square$

For verification, outcome states (faulted and repairable) have $\beta=V+L_f-c>0$. Cost-only states (clean,
unrepairable, or uninformative) have $\beta=-c$. Hence:

\begin{corollary}[vote rule]
GRPO verifies iff $P(\text{outcome state})>P(\text{cost-only state})$, i.e.\ $p\,s>1/2$, independent of $c$, $V$ and
$L_f$. Dr.~GRPO and RLOO verify iff $p\,s\,(V+L_f)>c$.
\end{corollary}

Nothing in the theorem is specific to reading. Any binary choice whose payoff is fixed by a hidden state shared within
the group is decided by a count of states, not by stakes: a rare catastrophic state casts the same vote as a common
trivial one.

\paragraph{When it weakens or vanishes.} The exact enumeration (Table~\ref{tab:thr}) shows six regimes:
\begin{itemize}
\item \textbf{Exogenous noise.} Within-state noise that is not caused by the decision (an exogenous failure rate
$q_0=0.3$) restores cost sensitivity, biased upward at small $c$. This is the counterexample: GRPO does \emph{not}
remove costs in general.
\item \textbf{Stochastic recovery} ($s<1$) shrinks the outcome states' vote. At $s=0.3$, GRPO never verifies.
\item \textbf{Independent states per rollout} (condition B) weaken but do not remove the bias, because the group
scale still depends on the sampled composition of states.
\item \textbf{Free reads} ($c=0$). Cost-only groups have zero variance, so GRPO verifies whenever $p>0$. This
matches the earlier free-read result (read-back 1.00 in five of five seeds).
\item \textbf{A never-reading policy} ($\rho=0$) is absorbing for every estimator. This matches the ID1 result: a
warm start that never reads stays at 0/153.
\item \textbf{Very small costs} ($c\,\sigma_r\lesssim\epsilon$): the $\epsilon$ term restores some dependence on the
cost.
\end{itemize}
The general condition is that \emph{a scale preserves every decision's expected sign iff it is independent of the
hidden states of the samples it scales}. No scale, RLOO, and an EMA of past scales qualify. Group std, small-batch
std, fork- or tree-node std and GDPO's per-component std \citep{liu2026gdpo} do not.

\begin{table}[t]
\centering\small
\begin{tabular}{lccc}
\toprule
Regime ($G{=}8$, $\rho{=}0.3$) & $c{=}0.01$ & $c{=}0.05$ & $c{=}0.2$\\
\midrule
GRPO, A, $s{=}1$ & 0.495 & 0.499 & 0.500\\
GRPO, A, $s{=}0.6$ & 0.654 & 0.667 & 0.706\\
GRPO, A, $q_0{=}0.3$ & 0.090 & 0.134 & 0.281\\
GRPO, B, $s{=}1$ & 0.233 & 0.247 & 0.310\\
Dr.~GRPO / RLOO, $s{=}1$ & 0.010 & 0.050 & 0.200\\
\bottomrule
\end{tabular}
\caption{Exact fault-rate threshold $p^*$ above which the expected update favours verifying
(\texttt{results/lab/exact\_analysis.json}). The VOI threshold is $c/(s(V{+}L_f))$; Dr.~GRPO and RLOO attain it
exactly.}
\label{tab:thr}
\end{table}

%%SECTION4%%

\section*{Limitations}

\begin{itemize}
\item Theorem~\ref{thm:vote} is exact for one binary decision with deterministic within-state payoffs. Multi-decision
trajectories, token-level loss weighting, KL and clipping are covered only empirically.
\item The testbed is tabular and its policy class is small, so every unbiased estimator eventually reaches the
optimum. Speed differences there do not imply different end points in an LLM.
\item The LLM evidence is from one model family and one simulator, with development seeds that produce identical
greedy behaviour. The confirmatory AD-C1 comparison has not run.
\item The judge study has one human annotator.
\item Every fork-based estimator needs a resettable simulator.
\end{itemize}

\section*{Ethics Statement}

Synthetic purchase orders only; no real payment, person or organisation is involved. The mechanism we describe can
make a trained agent skip checks whose stakes are high but whose failure states are rare. We recommend normalisers
that do not depend on the hidden state, together with deterministic verification of writes.

\bibliography{refs,refs_v2}
\bibliographystyle{acl_natbib}

\appendix
\section{Derivations}\label{app2:proofs}

\paragraph{Theorem~\ref{thm:vote}.} Fix a state $z$ with $R_i=\alpha_z+\beta_z r_i$. Then
$R_i-\bar R=\beta_z(r_i-\bar r)$ and $\hat\sigma_R=|\beta_z|\hat\sigma_r$, so
$\hat A_i=\sgn(\beta_z)(r_i-\bar r)/\hat\sigma_r$ whenever $0<k<G$; otherwise $\hat A_i=0$. The score of the verify
logit under a Bernoulli policy is $r_i-\rho$. Because $\sum_i(r_i-\bar r)=0$,
$\sum_i\hat A_i(r_i-\rho)=\sgn(\beta_z)\sum_i(r_i-\bar r)^2/\hat\sigma_r=\sgn(\beta_z)(G-1)\hat\sigma_r$.
Its expectation over $k\sim\mathrm{Bin}(G,\rho)$ is $(G-1)m_G(\rho)\sgn(\beta_z)$. Without the std the same steps give
$\beta_z\sum_i(r_i-\bar r)^2$, with expectation $\beta_z(G-1)\rho(1-\rho)$. RLOO's advantage is $G/(G-1)$ times the
centred one. For Dr.~GRPO and RLOO only i.i.d.\ rollouts given the group's state(s) are used:
$\mathbb E\sum_i(R_i-\bar R)(r_i-\rho)=(G-1)\mathrm{Cov}(R,r)=(G-1)\rho(1-\rho)(\mathbb E[R\mid r{=}1]-\mathbb E[R\mid r{=}0])$.
This holds for stochastic payoffs and under condition B. Averaging over $z$ completes the proof. All three
identities are checked by exact enumeration of multinomial group compositions
(\path{actionproof/lab/exact_vote.py}; 21 tests).

\paragraph{Normaliser condition.} If $s$ is independent of the hidden states and returns of the samples it scales,
given $x$, then $\mathbb E[\sum_iA_i\nabla\log\pi/s\mid x]=\mathbb E[1/s\mid x]\,\mathbb E[\sum_iA_i\nabla\log\pi\mid x]$.
This is a positive rescaling of each task's unbiased gradient, so every decision keeps its expected sign. Any $s$
computed from the same group violates the independence: in a clean group $s\approx c\,\hat\sigma_r$, while in a
faulted, repairable group $s\approx(V+L_f)\hat\sigma_r$.

\paragraph{Paired-fork decision value.} For an on-policy state $S_t$ with history $h_t$, forced action $a$,
continuation by the current $\pi$ on the branch's own history, and any coupling of exogenous noise that preserves
marginals,
$\mathbb E[G_a\mid S_t]=Q^\pi(h_t,z,a)$ and $\mathbb E_z[Q^\pi(h_t,z,a)\mid h_t]=Q^\pi(h_t,a)$. The paired contrast is
therefore unbiased for $\Delta^\pi(h_t)$ for any number of forks. CRN reduces only
$\mathbb E_z\mathrm{Var}(\hat\Delta\mid z,h)$; the spread $\mathrm{Var}_z(\Delta^\pi(h,z))$ remains. A stored fork mean
used $m$ updates later is biased by at most $2HG_{\max}\max_h\mathrm{TV}(\pi_k,\pi_{k+m})$ (performance-difference
lemma). Full statements are in \path{docs/research_lab/CDVO_THEORY.md}.

\section{Algorithms}\label{app2:alg}

\paragraph{Paired-fork credit (CDVO family).} At each step, for each selected verification decision $(h_t,S_t)$:
draw a seed $\xi$; for each action $a$ of the decision, restore $S_t$, force $a$, then sample the continuation from
$\pi_\theta$ on the branch's own history with exogenous draws from $\xi$, and record $G_a$. The gradient term is
$\sum_a\pi(a\mid h_t)\big(G_a-\sum_b\pi(b\mid h_t)G_b\big)\nabla\log\pi(a\mid h_t)$, and every unselected decision keeps the
sequence term $(R-\bar R_g)\nabla\log\pi$. Selection: \emph{full} (every decision), \emph{fixed} (each with
probability $b$), or \emph{adaptive} (the same expected count, given to decisions whose observable context has the
smallest $|\bar\Delta|/\mathrm{se}$ over its last 20 fork contrasts from earlier steps). Training stops at a fixed
episode budget, with forks charged as partial episodes.

\paragraph{Information-gated credit (igc\_v1).} A first read after a write whose observation is identical in every
state consistent with $h_t$ gets the exact all-action term with $Q=(0,-c)$. All other decisions use the centred
sequence term.

\section{Environment, data, licences and splits}\label{app2:data}

\paragraph{ProcureSim}\label{app2:env} (\path{actionproof/integrity/sim/}) exposes five read tools, six write tools
and a report tool. Every write passes one chokepoint, which applies the task's fault plan and records whether the
write committed nothing, part, all, or twice. The verifier reads only the final state. An independently written
reference verifier and 91 golden traces agree with it (\path{tests/golden/}). The costed pool \texttt{cw}
charges $c=0.05$ per read (\path{actionproof/ad/rewards.py}).

The testbed tasks are generated from fixed seeds (\path{actionproof/ad2/tabular/envs.py}). The training and
evaluation workflows are fixed; the generalisation splits (unseen fault rates, repair rates, costs and workflow
length) come from disjoint seeds. ProcureSim splits are generated by \path{actionproof/bench/generator_proc.py};
task and purchase-order ids are disjoint across splits, and the prompt never depends on the fault plan (both tested).
Hashes are in \path{results/*_hashes.json} and \path{docs/stage2/BENCHMARK_CARD.md}. All data are synthetic.

Models and licences:
\begin{itemize}
\item policy: Tiny Aya Global (Cohere; licence per its model card);
\item judges (Hub revisions in \path{results/j2_model_revisions.jsonl}): Gemma~4 31B (the Hub metadata reads
apache-2.0; J1 listed ``Gemma licence'', to be checked), Qwen3-235B-A22B-Instruct-2507 (Apache-2.0), gpt-oss-120b
(Apache-2.0), DeepSeek-V3.2 (MIT) and Aya Expanse 32B (CC-BY-NC-4.0).
\end{itemize}
Leakage checks: \path{tests/integrity/test_procsim.py} and the actor-input test
\path{tests/lab/test_cdvo.py::test_policy_inputs_never_contain_the_hidden_world}.

\section{Training configuration and compute}\label{app2:train}

\paragraph{Testbed.} The policy has one softmax head per decision type over [bias, op-type one-hot, known-failed]
(tabular) or 8 observable descriptors (feature policy). Settings shared by every estimator: Adam with lr 0.05; 4
groups of 8 rollouts per step; warm start $\rho_0=0.3$ and repair 0.9; loss summed over decisions and divided by
their number (TRL `dapo').

\paragraph{PPO} uses a tabular EMA critic, 4 epochs and clip 0.2.

\paragraph{Budgets} are E1 = 6{,}400 and E4 = 25{,}600 episodes. The exact optimum and policy values are computed by
expectimax.

\paragraph{Compute.} All testbed runs were on Slurm CPU nodes (\texttt{compute}, 32 cores; job 467649). An earlier
local attempt was stopped and is not used.

\paragraph{LLM.} TRL GRPO: 200 steps; 4 prompts $\times$ 8 rollouts; lr $2\times10^{-5}$; KL 0.02; LoRA rank 16; one
H200, H100 or A100 per run (about 1.5--3.5\,h). Every run is in \path{results/run_manifest.json}, which records
commit, configuration, adapter and data hashes, seed, steps, status, and the CPU and judge runs.

\section{Per-seed and statistical tables}\label{app2:results}

\paragraph{Registration.} LAB1 (\texttt{PREREGISTRATION\_LAB1.md}) was frozen before any registered output; its
dry-run and its seed-98/99 smoke test are disclosed there. The primary rule required, at E4: $\geq20\%$ lower
seed-mean regret than the strongest simple baseline (SFB), better in $\geq4/5$ seeds, never-faulted verification and
duplicate loss within tolerance, and Holm-adjusted one-sided paired sign-flip $p<0.05$ over 60 (setting, seed)
units. Table~\ref{tab:gates} gives the outcome; per-seed regrets are in Table~\ref{tab:perseed}.

\begin{table}[h]\centering\small\begin{tabular}{lcccccc}
\toprule
Test & SFB & cand. & SFB & rel. & seeds & Holm $p$\\
\midrule
H1 E1 & ppo & 0.0595 & 0.0215 & -1.76 & 0/5 & 1.00 \\
H1 E4 & nostd & 0.0170 & 0.0080 & -1.14 & 0/5 & 1.00 \\
H2 E1 & igc\_v1 & 0.0380 & 0.0202 & -0.88 & 0/5 & 1.00 \\
H2 E4 & igc\_v1 & 0.0097 & 0.0062 & -0.57 & 0/5 & 1.00 \\
\bottomrule
\end{tabular}

\caption{LAB1 primary tests. H1: adaptive forks vs SFB over \{Dr.~GRPO, RLOO, value baseline, learned Q, PPO\}.
H2: adaptive forks + igc vs SFB including igc. Neither is supported.}\label{tab:gates}\end{table}

\begin{table}[h]\centering\resizebox{\columnwidth}{!}{\begin{tabular}{lcccccc}
\toprule
Estimator & s0 & s1 & s2 & s3 & s4 & mean\\
\midrule
Dr.~GRPO & 0.0075 & 0.0078 & 0.0086 & 0.0080 & 0.0079 & 0.0080 \\
RLOO & 0.0075 & 0.0078 & 0.0081 & 0.0084 & 0.0080 & 0.0080 \\
PPO & 0.0138 & 0.0089 & 0.0065 & 0.0141 & 0.0082 & 0.0103 \\
Value baseline & 0.0114 & 0.0085 & 0.0103 & 0.0111 & 0.0104 & 0.0103 \\
Learned Q & 0.0219 & 0.0080 & 0.0062 & 0.0135 & 0.0156 & 0.0131 \\
igc & 0.0060 & 0.0057 & 0.0064 & 0.0060 & 0.0070 & 0.0062 \\
AWRB-n & 0.0079 & 0.0081 & 0.0072 & 0.0072 & 0.0073 & 0.0075 \\
Forks, full & 0.0210 & 0.0194 & 0.0226 & 0.0227 & 0.0224 & 0.0216 \\
Forks, full, no CRN & 0.0220 & 0.0210 & 0.0250 & 0.0218 & 0.0220 & 0.0224 \\
Forks, random 25\% & 0.0149 & 0.0131 & 0.0147 & 0.0151 & 0.0138 & 0.0143 \\
Forks, adaptive 25\% & 0.0166 & 0.0170 & 0.0167 & 0.0168 & 0.0181 & 0.0170 \\
Forks, adaptive 10\% & 0.0122 & 0.0131 & 0.0147 & 0.0133 & 0.0118 & 0.0130 \\
Adaptive + stored-mean critic & 0.0911 & 0.0808 & 0.0731 & 0.0946 & 0.0859 & 0.0851 \\
Adaptive + igc & 0.0106 & 0.0087 & 0.0097 & 0.0105 & 0.0092 & 0.0097 \\
\bottomrule
\end{tabular}
}
\caption{Exact regret per seed at E4 (25{,}600 episodes; mean over 12 settings $\times$ 3 workflows).}
\label{tab:perseed}\end{table}

\paragraph{Normaliser intervention (L2N).} The grid is $p\in\{0.1,0.2,0.4,0.6\}$, $c\in\{0.01,0.05,0.2\}$,
$s\in\{0.6,1\}$, with 3 seeds. The wrong-decision rate counts disagreements with the sign of the exact VOI where
$|$VOI$|>0.02$ (Table~\ref{tab:l2n}).

\begin{table}[h]\centering\small\begin{tabular}{lcc}
\toprule
Scale & wrong-decision rate & regret\\
\midrule
grpo & 0.727 & 0.1249 \\
indep\_world\_std & 0.667 & 0.0888 \\
indep\_task\_std & 0.061 & 0.0129 \\
nostd & 0.045 & 0.0095 \\
\bottomrule
\end{tabular}

\caption{Same estimator, different scale: from 16 independent rollouts in the same hidden state
(\texttt{indep\_world\_std}) or in fresh states (\texttt{indep\_task\_std}). Registered H4 is
supported.}\label{tab:l2n}\end{table}

\paragraph{Exact gradient study (L2G).} We draw 300 independent batches at the warm start and compare them with the
exact gradient obtained by finite differences of the exact value (Table~\ref{tab:grad}). GRPO's SNR is that of a
wrong-signed mean. Registered H3 (CRN lowers variance on every risky logit) is not supported: the variance ratios
are 0.94--1.12.

\begin{table}[h]\centering\small\resizebox{\columnwidth}{!}{\begin{tabular}{lccccc}
\toprule
Estimator & cos & order & receive & pay & never-faulted\\
\midrule
GRPO & -0.175 & 0.97 & 0.84 & 0.87 & 1.80 \\
Dr.~GRPO & +0.995 & 0.12 & 0.35 & 0.33 & 0.34 \\
igc & +0.998 & 0.12 & 0.35 & 0.36 & 23.48 \\
Forks & +1.000 & 0.14 & 0.38 & 0.46 & 12.43 \\
Forks, no CRN & +0.996 & 0.19 & 0.42 & 0.43 & 13.03 \\
\bottomrule
\end{tabular}
}
\caption{Cosine of the mean gradient with the exact gradient, and per-logit SNR $|\text{mean}|/\text{sd}$ on the
verify logits (rich multi, $c=0.05$, $L_f=1$).}\label{tab:grad}\end{table}

\paragraph{Recovery-start agent (S3, confirmatory).} Table~\ref{tab:s3v2} is reproduced from the registered S3
analysis (\path{results/s3/s3_summary.json}, \path{s3_addendum2.json}).
\begin{table}[h]\centering\resizebox{\columnwidth}{!}{% generated by scripts/paper_tables.py; do not edit
\begin{tabular}{llccc}
\toprule
Contrast & Split & $\Delta$VSC & $\Delta$violation & Seeds $>0$ (VSC) \\
\midrule
RSC $-$ R$_V$ (recipe) & recovery starts & +0.21 [+0.17, +0.25] & -0.11 [-0.15, -0.08] & 5/5 \\
RSC $-$ R$_V$ (recipe) & full episodes & +0.21 [+0.17, +0.25] & -0.12 [-0.15, -0.08] & 5/5 \\
RSC $-$ R$_V$ (recipe) & stress & +0.39 [+0.30, +0.48] & -0.39 [-0.48, -0.30] & 5/5 \\
RSC $-$ matched (start state) & recovery starts & +0.09 [+0.06, +0.11] & +0.01 [-0.00, +0.02] & 3/3 \\
RSC $-$ matched (start state) & partial + duplicate & +0.22 [+0.15, +0.29] & +0.03 [+0.00, +0.05] & 3/3 \\
RSC $-$ matched (start state) & stress & -0.04 [-0.07, -0.01] & +0.04 [+0.01, +0.07] & 1/3 \\
RSC $-$ matched, 5 seeds (Holm $p$=0.036) & recovery starts & +0.11 [+0.08, +0.13] & -- & 5/5 \\
RSC $-$ matched, 5 seeds (Holm $p$=0.111) & partial (full episodes) & +0.38 [+0.32, +0.45] & -- & 4/5 \\
RSC $-$ matched, 5 seeds (Holm $p$=0.337) & duplicate (full episodes) & +0.13 [+0.09, +0.17] & -- & 2/5 \\
\bottomrule
\end{tabular}
}
\caption{Registered S3 contrasts: seed-averaged difference [95\% task CI].}\label{tab:s3v2}\end{table}

\section{Ablations, generalisation and negative results}\label{app2:neg}

\begin{table}[h]\centering\small\begin{tabular}{lcccc}
\toprule
Estimator & train & iid & shift & long\\
\midrule
Dr.~GRPO & 0.0074 & 0.0058 & 0.0045 & 0.0295 \\
RLOO & 0.0074 & 0.0058 & 0.0045 & 0.0295 \\
PPO & 0.0019 & 0.0024 & 0.0028 & 0.0243 \\
Learned Q & 0.0031 & 0.0039 & 0.0038 & 0.0291 \\
igc & 0.0044 & 0.0046 & 0.0045 & 0.0362 \\
Forks, full & 0.0185 & 0.0193 & 0.0137 & 0.0781 \\
Forks, adaptive & 0.0096 & 0.0103 & 0.0089 & 0.0574 \\
Adaptive + igc & 0.0090 & 0.0078 & 0.0065 & 0.0474 \\
\bottomrule
\end{tabular}

\caption{Generalisation (feature policy, 51{,}200 episodes, 5 seeds): exact regret on training tasks and on
unseen splits (iid; shifted $p$, $s$, $c$; longer workflows).}\label{tab:l3g}\end{table}

\begin{table}[h]\centering\resizebox{\columnwidth}{!}{\begin{tabular}{lcccccc}
\toprule
Run & 0--20 & 20--50 & 50--100 & 100--150 & 150--200 & zero-var\\
\midrule
ad\_grpo\_s0 & 0.19 & 0.04 & 0.04 & 0.04 & 0.02 & 0.81 \\
ad\_grpo\_s1 & 0.20 & 0.03 & 0.04 & 0.03 & 0.04 & 0.66 \\
ad\_nostd\_s0 & 0.23 & 0.19 & 0.26 & 0.38 & 0.48 & 0.03 \\
ad\_nostd\_s1 & 0.28 & 0.26 & 0.28 & 0.42 & 0.50 & 0.05 \\
ad\_awrb\_s0 & 0.30 & 0.61 & 0.77 & 0.74 & 0.76 & 0.14 \\
ad\_awrbn\_s0 & 0.25 & 0.23 & 0.28 & 0.43 & 0.45 & 0.09 \\
\bottomrule
\end{tabular}
}
\caption{Language agent, costed pool (development): share of sampled training rollouts that read, by step window,
and the share of zero-variance groups in steps 150--200 (\path{results/lab/train_read_rate.json}).}
\label{tab:rr}\end{table}

\paragraph{Negative results kept.}
\begin{itemize}
\item Forked counterfactual credit in the LLM (FCA; $-0.10$ VSC; dropped by its kill rule).
\item Paired-fork credit in the testbed at matched budget (all variants, 0/5 seeds); adaptive allocation no better
than random; a stored-fork-mean critic that collapses verification (regret 0.085).
\item CRN without variance reduction.
\item AWRB, whose gain is over-verification and which is fragile to prior misspecification.
\item Per-context EMA scaling and the EMA-scaled LLM arms (one seed of each collapsed to 0 reads).
\item Condition-B GRPO still biased (0.315).
\item One third of igc's earlier gain was a testbed artefact.
\item Turn-credit (C5) and hint-guided sampling (G), dropped.
\item Tool-name randomisation (SR1), not supported.
\item Negative transfer to a sibling domain.
\end{itemize}
Each is documented with a diagnosis in \path{docs/research/} and \path{docs/research_lab/LAB1_RESULTS.md}.

\section{LLM judge validation}\label{app2:judge}

%%JUDGE%%

\section{Representative trajectory}\label{app2:traj}

Task \texttt{s3\_conf\_full-0014} (\path{results/s3/conf/}): every model pays 62{,}892 of 251{,}570 cents and
reports ``completed''. The trained models read the payment record first, so reading did not lead to recovery. In the
costed pool, the GRPO arm's episodes contain no read at all
(\texttt{verify\_after}~=~0 for every operation, \path{results/ad/d1_summary.json}).

\section{Reproduction and declarations}\label{app2:repro}

\paragraph{Reproduction.}
\begin{itemize}
\item \path{python scripts/lab/exact_analysis.py} reproduces Table~\ref{tab:thr} and the TB1 recount.
\item \path{sbatch slurm/lab/lab1_cpu.sbatch} reproduces LAB1, gated by \path{scripts/lab/freeze_lab1.py
--check}.
\item \path{python scripts/ad2/decompose.py ...} reproduces the log decomposition.
\item \path{python evaluation/llm_judge_committee/analyze_j2.py} reproduces the judge tables.
\item \path{scripts/s2_numcheck.py} checks every number against \path{claims/claims.yaml}.
\end{itemize}

\paragraph{Kits and models.} AgentTune provides the batched GRPO rollouts and AuditKIT the Stage~1 judge audit; the
models are Cohere's.

\paragraph{Safety and general capability.} Registered IFEval and BFCL v4 retention (R1): GRPO is non-inferior to its
warm start within 0.03.

\paragraph{ROI.} The fix that matters, dropping the group std, costs nothing at training time. The forked estimators
spend 1.9--5.1$\times$ the episodes per policy update, for no gain at matched budget.

\paragraph{Prior work and team.} ActionProof Stage~1 (Nov 1--15) built the e-commerce environment. Stage~2 (Nov
20--27) built ProcureSim, the costed pool, the testbed, the judge committee and every result here. Team
contributions: \pend{author to state}.

\paragraph{LLM use.} Claude (Anthropic) was used as a coding, analysis and writing assistant. Every number was
produced by code in the repository and is checked against its result file.


\end{document}
